\begin{Huge}
    Kognitionswissenschaft-FAQ
\end{Huge}

\begin{block}{H\"aufig gestellte Fragen zum Studium}

\question{Lernt man im Studium, wie man programmiert?}{
    Ja, aber auf einer sehr eigenst\"andigen Basis. Man bekommt einen \"Uberblick \"uber die Sprache(n) und \"uber dahinterstehende Programmierkonzepte. Praktische Programmiererfahrung muss man sich aber gr\"o"stenteils selbst erarbeiten.
}

\question{Welche Programmiersprachen macht man da so?}{
    Das ist von Prof zu Prof und von Veranstaltung zu Veranstaltung unterschiedlich. In den ersten beiden Semestern lernt man meistens Racket und Java.
}

\question{Muss man programmieren k\"onnen, um das Studium anzufangen?}{
    Nein. Die Vorlesung beginnt absolut bei 0, um allen den Einstieg zu erm\"oglichen.
}

\question{Muss man gut in Mathe sein?}{
    Man muss kein Mathe-Genie sein, man sollte Mathe aber nicht hassen. Es ist gerade am Anfang viel Mathe.
}

\question{Was mache ich nachher damit?}{
    Eine wirklich sehr gute Frage - die nicht so einfach zu beantworten ist. Wir haben die M\"oglichkeit uns im Master weiter zu spezialisieren, Richtung Informatik, Biologie oder was sonst noch so grob passt. Davon h\"angt dann auch sehr stark die sp\"atere Berufsaussicht ab. Reine Kognitionswissenschaft findet man vor allem in der Forschung, ein weiteres gro"ses Feld ist Maschinelles Lernen, bzw. K\"unstliche Intelligenz.
}

\question{Wie ist die Frauenquote so?}{
    58,8\%.
}

\question{Ich kenne mich mit Informatik gar nicht aus. Ist das schlimm?}{
    Nein, ist es nicht. Wenn du einigerma"sen logisch denken kannst und dir auch Mathe ganz gut liegt, dann kriegst du das auch hin.
}

\question{Eigentlich will ich ja Psychologie studieren, aber mein NC reicht nicht. Kogni ist doch auch was mit Psychologie, oder?}{
    Schon, ja. Allerdings ist der Anteil im Bachelor-Studiengang nicht allzu hoch. Besonders am Anfang sind gr\"o"sere Schwerpunkte Informatik und Mathe, das sollte man auf jeden Fall mit in Betracht ziehen. Im weiteren Verlauf des Studiums kann man sich aber weiter in Richtung Psychologie vertiefen.
}

\question{Kann ich danach auf Psychologie wechseln?}{
    Das k\"onnen wir so allgemein nicht beantworten. Generell h\"angt das sehr stark von der Uni ab, an der man studieren m\"ochte. Das m\"usste man dann im Zweifelsfall vorher schon mal anfragen.
}

\question{Gibt es Praktika?}{
    Im normalen Studienverlauf ist kein berufsorientiertes Praktikum vorgesehen, viele arbeiten aber parallel als Werksstudent*in oder man macht ein Kurzpraktikum in den Semesterferien.
}

\question{Kann man ein Auslandssemster machen?}{
    Klar, geht immer. T\"ubingen nimmt am ERASMUS-Programm teil, die Organisation ist aber langwierig und man sollte sich fr\"uh (ein Jahr vorher) drum k\"ummern.
}

\question{Wie ist da so der NC?}{
    1,9 (WS26/27). Das muss aber in den n\"achsten Jahren nicht zwingend noch so sein, der Wert \"andert sich immer.
}
\end{block}
